% !Mode:: "TeX:UTF-8"
\chapter{绪论}

\section{课题来源与背景}

迭代学习控制（Iterative Learning Control, ILC）的概念由 Arimoto 等人于 1984 年首次提出
\upcite{arimoto1984}，其核心思想是：对于在有限时间区间内重复执行同一跟踪任务的系统，利用
前一次迭代的误差信息修正当前次的控制输入，从而沿迭代轴逐步提升跟踪精度。ILC 的提出源于
工业机器人重复操作中对高精度轨迹跟踪的迫切需求——传统反馈控制因不具备学习与记忆机制，
在每次重复执行时出现相同的跟踪误差，无法利用历次运行积累的经验改善控制品质\upcite{moore1993}。
ILC 恰好弥补了这一不足，为重复运行系统开辟了"从迭代中学习"的控制范式\upcite{longman2000}。在理论上，ILC
经历了从早期 D 型、P 型学习律的频域收敛分析\upcite{moore1992}，到基于状态空间模型的时域
设计\upcite{saab1995}，再到以范数优化为核心的凸优化框架\upcite{amann1996,gunnarsson2001,owens2005}
的发展路径，逐步构建起从启发式设计到系统化优化理论的完备体系\upcite{ahn2007,bristow2006,longman2000,xu2011survey}。
在应用上，ILC 已成功部署于工业机器人\upcite{norrlof2002}、精密运动系统\upcite{bolder2015}、
半导体晶圆台\upcite{mishra2009}、增材制造\upcite{wang2018}及康复医疗\upcite{freeman2012}等
广泛领域，展示出显著的工程价值。然而，ILC 的性能高度依赖于被控系统的精确动态模型或充分
激励的输入输出数据，如何在有限数据条件下保证学习性能，是该领域长期面临的核心挑战\upcite{owens2016}。

在迭代学习控制的理论框架中，范数最优迭代学习控制（Norm-Optimal ILC, NOILC）由 Amann
等人 \upcite{amann1996} 提出，其将控制输入更新问题转化为一个凸优化问题，在 Hilbert 空间
框架下最小化跟踪误差与输入增量的加权范数之和。范数最优迭代学习控制无需手动调节参数即可
保证跟踪误差范数沿迭代轴单调收敛，因而被广泛视为基于模型方法的性能标杆。然而，范数最优
迭代学习控制的设计与实现高度依赖被控系统的解析模型——具体而言，需要已知由 Markov 参数
构成的 Toeplitz 矩阵 $\mathbf{G}$ 以及系统初始响应 $\mathbf{d}$。在实际工程中，复杂系统
的精确数学模型往往难以获取\upcite{hou2013}：一方面，许多物理过程（如柔性结构、软体机器人
等）机理复杂、参数众多，建立高保真模型的成本高昂且精度有限；另一方面，系统特性可能随
工况、老化和环境变化而漂移，基于固定模型设计的控制器缺乏对模型失配的鲁棒性。这一矛盾
构成了基于模型迭代学习控制方法在实际应用中的根本瓶颈。

为摆脱对精确数学模型的依赖，数据驱动控制（Data-Driven Control）成为近年来控制领域的
研究热点，其基本理念是直接利用从系统采集的输入输出数据设计控制器，绕过显式的系统辨识
环节\upcite{ljung1999,fliess2013}。在这一方向中，Willems 等人发展的行为系统理论 \upcite{willems1991} ，以及基于行为系统理论提出的 Willems' 基本引理 \upcite{willems2005} 提供了重要的理论基础。Willems' 基本引理指出：若离线采集的
输入数据满足持续激励（Persistency of Excitation, PE）条件，则任一条由同一系统生成的长度
为 $L$ 的输入输出轨迹均可表示为离线输入输出数据所构造的联合 Hankel 矩阵的列的线性组合。换言之，
一段信息量足够丰富的离线数据可以完整刻画系统的全部动态行为，从而在数学上等价于系统的
解析模型。基于这一引理，Markovsky 和 Rapisarda \upcite{markovsky2008} 提出了数据驱动仿真
与控制的基本框架，为后续数据驱动迭代学习控制的研究奠定了基础。

围绕 Willems' 基本引理，数据驱动控制研究在随后的十余年间取得了丰硕成果。De Persis 和 Tesi
\upcite{depersis2019,depersis2020} 利用该引理导出了数据驱动状态反馈镇定与线性二次型调节器（LQR）的直接
设计公式，首次在理论上建立了数据驱动控制与模型基控制之间的等价关系。Van Waarde 等人
\upcite{vanwaarde2020} 将基本引理从单数据集推广至多数据集情形，并进一步提出了"数据信息量"（Data
Informativity）的统一分析框架——从数据是否包含充分系统动态信息的角度，重新诠释了数据驱动
分析与控制的数学基础 \upcite{vanwaarde2023}。Coulson、Lygeros 与 Dörfler \upcite{coulson2019}
提出了数据赋能预测控制（Data-Enabled Predictive Control, DeePC），将传统模型预测控制的
滚动时域优化问题改造为仅依赖 Hankel 数据矩阵的等价形式，在电力系统与建筑节能等应用中展示
了良好的控制性能。Berberich 等人 \upcite{berberich2021,berberich2022} 在此基础上进一步发展了
具有稳定性与鲁棒性保证的数据驱动模型预测控制框架，利用噪声数据的矩阵扰动分析建立了
闭环稳定性理论。Markovsky 与 Dörfler \upcite{markovsky2021} 系统综述了行为系统理论在数据驱动
分析、信号处理与控制中的各类应用，梳理了从 Willems' 基本引理到数据驱动控制器设计的完整
技术脉络；Dörfler 等人 \upcite{dorfler2022} 进一步阐明了直接与间接数据驱动控制
 formulation 之间的内在联系。上述工作充分表明，Willems' 基本引理已成为连接行为系统理论与现代数据驱动控制的
核心桥梁，为"绕过系统辨识、直接由数据到控制器"的设计范式提供了坚实的理论支撑。

近年来，已有学者将数据驱动方法引入迭代学习控制领域，提出了数据驱动的范数最优迭代学习
控制框架\upcite{janssens2013}，并在理论上证明了其与基于模型的范数最优迭代学习控制完全等价——即在不使用系统
解析模型的前提下，仅凭输入输出数据即可实现相同的跟踪性能 \upcite{jiang2025}。然而，这一
等价性建立在输入数据满足持续激励条件的基础之上。具体而言，考虑 $m$ 维输入、$n$ 维状态的
离散线性时不变系统，若需利用长度为 $J$ 的离线输入数据表示长度为 $L$ 的轨迹，则 离线输入数据须满足阶数为 $L+n$ 的持续激励
条件，这要求 $J\geq (m+1)(L+n)-1$。当系统阶数较高或控制时域较长时，所需数据量可能极为庞大。过长的离线测试不仅显著增加实验时间与经济成本，还可能
因长时间施加测试信号引发安全隐患（如激励系统的不稳定模态或超出执行器物理限幅）
\upcite{vanwaarde2020}，并加重后续在线控制计算的负担。

针对持续激励条件导致数据采集成本过高的问题，Camlibel 等人 \upcite{camlibel2025} 在系统
辨识的研究中提出了一种最短测试实验设计方法。该研究指出，只要新施加的输入避开了某个低维
仿射集，由输入输出数据构造的联合 Hankel 矩阵的秩便可严格增加 1，直至达到理论最大值，
从而以理论最短的数据长度满足 Willems' 基本引理的应用条件。本文借鉴了这一工作中的秩增长
引理作为数据生成工具，但研究目标与之不同：Camlibel 等人关注的是线性系统的辨识问题，而
本文致力于将最短数据生成机制与迭代学习控制的设计有机融合，构建一套从离线数据采集到在线
迭代更新的完整数据驱动迭代学习控制方案，并对其收敛性给出严格的理论保证。本文正是在此
背景下展开研究。

\section{选题意义}

在现代工业生产与精密制造中，大量系统需要在有限时间区间内重复执行同一跟踪任务，例如
工业机械臂的拾取与放置操作、数控机床的周期性轮廓加工、半导体晶圆台的往复扫描运动等。
对于这类重复运行系统，传统的反馈控制（如 PID 控制、最优控制）在每次执行时以相同的
控制器参数运行，不会利用历次运行中积累的误差信息来改善控制性能，上一次任务中出现
的跟踪误差在下一次执行时仍会以相同的方式复现。迭代学习控制正是针对这一局限而提出的
\upcite{bien1998}：它利用系统的重复运行特性，沿迭代轴从过去的误差中学习，逐步修正
控制输入，从而实现对期望轨迹的渐近精确跟踪\upcite{wang2009survey}。

基于模型的范数最优迭代学习控制虽然具有单调收敛等优良性质，但其设计依赖系统的精确数学
模型，而这在实际工程中往往难以获取\upcite{yin2015}；基于 Willems' 基本引理的数据驱动框架
可以绕开显式建模\upcite{polderman1998}，却要求输入数据满足持续激励条件，导致离线数据采集量
过大。两者之间尚缺少一种方案，能够以理论最短的数据长度获取信息量充分的离线数据，并将其
嵌入迭代学习控制的优化与控制更新之中。探索这一问题，有助于降低数据驱动迭代学习控制的
应用门槛，也可为最少数据与最优控制性能之间的权衡提供新的理论认识。

本文的选题意义主要在三个方面。在方法上，本文将 Camlibel 等人针对系统辨识提出的秩增长
引理\upcite{camlibel2025} 引入迭代学习控制的离线数据生成阶段。该引理表明，每一步只需避开
一个低维仿射集，联合 Hankel 矩阵的秩便可严格增加 1，直至达到理论最大值。本文据此设计了
最短测试实验方案，使离线数据长度达到理论最小值，远低于传统持续激励条件所要求的数据量，
显著压缩了离线实验的时间与经济成本。在理论上，本文将上述最短数据生成机制与范数最优迭代
学习控制的凸优化框架相结合，导出了仅依赖数据矩阵的控制更新律解析解，并证明了只要联合
Hankel 矩阵达到目标秩，算法便能保证跟踪误差范数单调收敛至零，且在伴随系统正定时具有
指数收敛速率，与基于模型的范数最优迭代学习控制完全一致。在应用上，本文方法构成了一套
完整的数据驱动迭代学习控制方案，从离线实验设计、数据采集截止判断到在线迭代控制更新，
全部流程仅依赖系统的输入输出数据，无需任何解析模型参数，对于模型结构未知或辨识成本高
的工程场景具有实际参考价值。

\section{研究内容与目标}

本文的研究围绕一个核心问题展开：对于模型未知的离散线性时不变系统，能否以理论最短的
离线数据长度，实现与基于模型的范数最优迭代学习控制相同的跟踪性能？围绕这一问题，本文
从离线数据生成、控制算法设计、收敛性分析和仿真验证四个环节逐次推进。

（1）最短测试实验数据生成。针对传统持续激励条件要求过长数据序列的问题，设计一种逐次避让低维仿射集的输入选取策略，使联合 Hankel 矩阵的秩
在每一步严格增加 1，直至达到理论最大值 $L m + n$。目标是以理论最短的数据长度
$(m+1)L + n - 1$ 获得信息量充分的离线输入输出数据，这与传统持续激励条件所需的离线输入数据长度相比节省了 $mn$ 个数据点，为后续控制设计奠定基础。

（2）基于最少输入输出数据的迭代学习控制算法设计。在最短实验数据的基础上，利用 Willems'
基本引理将系统行为表示为离线数据的线性组合，构建Markov参数矩阵的等价表示，进而将数据驱动的
范数最优迭代学习控制问题转化为仅依赖数据矩阵的凸优化问题，导出控制更新律的解析解。

（3）基于数据的迭代学习控制收敛性分析。对所提算法造成的跟踪误差进行严格的理论分析，证明在联合 Hankel 矩阵达到目标秩的
条件下，跟踪误差范数沿迭代轴单调收敛至零，且在伴随系统满足正定性条件时具有指数收敛
速率，与基于模型的范数最优迭代学习控制具有相同的收敛性质。

（4）数值仿真验证。选取具体的离散线性时不变系统作为仿真对象，首先验证最短测试实验设计
中联合 Hankel 矩阵秩的逐步增长过程，确认数据长度达到理论最小值；然后对不同复杂程度的
期望轨迹执行迭代学习控制，考察跟踪误差的收敛行为与最终跟踪精度，验证所提方法的有效性。

\section{论文章节安排}

本文围绕"以最少输入输出数据实现数据驱动迭代学习控制"这一主线，按照从理论基础到方法
设计再到仿真验证的顺序组织全文，共分为五章。

第一章为绪论，介绍迭代学习控制的研究背景与发展脉络，梳理范数最优迭代学习控制、数据
驱动控制与 Willems' 基本引理、最短测试实验等相关方向的已有进展，指出尚待解决的问题，
并阐述本文的选题意义、研究内容与目标。

第二章为预备知识与理论基础。首先给出离散线性时不变系统的状态空间描述及提升形式表示，
定义由 Markov 参数构成的 Toeplitz 矩阵 $\mathbf{G}$与加权范数等基本概念，建立迭代学习控制
的一般框架；然后系统回顾基于模型的范数最优迭代学习控制，介绍其优化问题表述、控制更新律
及收敛性质，作为后续数据驱动方法的性能基准。

第三章为数据驱动范数最优迭代学习控制。给出了最短测试实验引理及其证明，并且在行为系统理论的框架下，引入 Willems' 基本引理
及数据驱动仿真算法，给出零初始条件响应矩阵的构造方法，推导数据驱动的范数最优迭代学习
控制更新律，并证明其与基于模型的范数最优迭代学习控制的等价性。

第四章为基于最少输入输出数据的迭代学习控制。在第三章最短数据生成机制的基础
上，分析联合 Hankel 矩阵的秩条件与理论上界，阐明目标秩 $Lm+n$ 与数据信息量之间的
关系；然后给出完整的基于最少输入输出数据的范数最优迭代学习控制算法流程；最后对所提算法
进行严格的收敛性分析，证明跟踪误差范数的单调收敛性与指数收敛速率。

第五章为仿真验证。通过数值实验验证所提方法的有效性，包括联合 Hankel 矩阵秩的逐步增长
过程与理论最短数据长度的达成，以及在不同期望轨迹下的迭代学习控制跟踪性能。

最后为结论，总结本文的主要工作，并讨论有待进一步研究的问题。